% ============================================================
% Section 42: 一维单原子链的晶格波能量与熔点频率
% ============================================================
\newpage
\section{固体理论：一维单原子链的晶格波能量与熔点频率}
\label{sec:1d-chain-energy}

\begin{modelbox}[题目：晶格波的时间平均能量与熔点频率]
考虑一个纵波 $x_n=A\cos(\omega t-qna)$ 在原子质量为 $m$、间距为 $a$ 的单原子线型点阵中传播。只考虑最近邻相互作用，恢复力系数为 $\beta$。
\begin{enumerate}
    \item 求每个原子的时间平均总能量；
    \item 设一维晶体在熔点 $T_m$ 时的原子振动振幅是平衡间距的 $10\%$，证明熔点附近原子的振动频率
    \[
    \omega = \frac{10}{a}\sqrt{\frac{2k_BT_m}{m}}.
    \]
\end{enumerate}
\end{modelbox}

\begin{tipbox}[考点快速分析]
\begin{itemize}
    \item \textbf{动能与势能的时间平均}：谐振子中 $\langle K\rangle=\langle U\rangle$，总能量 $E=2\langle K\rangle$
    \item \textbf{键能的归属}：每根键由两个原子共享，计算单原子势能时要乘 $1/2$
    \item \textbf{色散关系的化简}：利用 $\omega^2=(4\beta/m)\sin^2(qa/2)$ 把势能表达式中的 $\beta$ 消去
    \item \textbf{能量均分定理}：经典极限下每个自由度的平均能量为 $k_BT$
    \item \textbf{Lindemann 判据}：熔点振幅 $\sim(0.1\text{--}0.15)a$ 是经验规律，本题取 $10\%$
\end{itemize}
\end{tipbox}

\begin{derivationbox}[标准解题过程]
\textbf{Step 1: 时间平均动能}

第 $n$ 个原子的瞬时动能
\begin{equation}
K_n = \frac12 m\dot x_n^{\,2} = \frac12 m\omega^2 A^2\sin^2(\omega t-qna).
\end{equation}
对周期 $T=2\pi/\omega$ 求时间平均，利用 $\langle\sin^2\rangle=1/2$：
\begin{equation}
\boxed{\langle K\rangle = \frac14 m\omega^2 A^2.}
\label{eq:1d-K-avg}
\end{equation}

\textbf{Step 2: 时间平均势能}

只计最近邻，第 $n$ 与 $n+1$ 原子间的键能为
\begin{equation}
U_{n,n+1} = \frac12\beta(x_n-x_{n+1})^2.
\end{equation}
每根键由两端原子共享，故第 $n$ 个原子分摊的势能为
\begin{equation}
U_n = \frac12 U_{n,n+1} = \frac14\beta(x_n-x_{n+1})^2.
\end{equation}

计算位移差（和差化积）：
\begin{align}
x_n-x_{n+1} &= A\bigl[\cos(\omega t-qna)-\cos(\omega t-q(n+1)a)\bigr] \nonumber\\
&= 2A\sin\frac{qa}{2}\,\sin\Bigl(\omega t-qna-\frac{qa}{2}\Bigr).
\end{align}
平方后对时间取平均，$\langle\sin^2\rangle=1/2$：
\begin{equation}
\overline{(x_n-x_{n+1})^2} = 4A^2\sin^2\frac{qa}{2}\cdot\frac12 = 2A^2\sin^2\frac{qa}{2}.
\end{equation}
因此
\begin{equation}
\langle U\rangle = \frac14\beta\cdot 2A^2\sin^2\frac{qa}{2} = \frac12\beta A^2\sin^2\frac{qa}{2}.
\label{eq:1d-U-avg-raw}
\end{equation}

\textbf{Step 3: 利用色散关系化简势能}

一维单原子链的色散关系为
\begin{equation}
\omega^2 = \frac{2\beta}{m}(1-\cos qa) = \frac{4\beta}{m}\sin^2\frac{qa}{2},
\end{equation}
即
\begin{equation}
\beta\sin^2\frac{qa}{2} = \frac14 m\omega^2.
\end{equation}
代入式 \eqref{eq:1d-U-avg-raw}：
\begin{equation}
\langle U\rangle = \frac12\cdot\frac14 m\omega^2 A^2 = \frac14 m\omega^2 A^2.
\end{equation}
与式 \eqref{eq:1d-K-avg} 比较，得到\textbf{能量均分}：
\begin{equation}
\boxed{\langle K\rangle = \langle U\rangle = \frac14 m\omega^2 A^2.}
\end{equation}

\textbf{Step 4: 每个原子的时间平均总能量}
\begin{equation}
\boxed{\langle E\rangle = \langle K\rangle + \langle U\rangle = \frac12 m\omega^2 A^2.}
\label{eq:1d-E-avg}
\end{equation}

\textbf{Step 5: 熔点频率的推导}

在经典极限下（熔点附近通常满足），一维谐振子每个自由度的平均能量为 $k_BT$（动能贡献 $\frac12k_BT$，势能贡献 $\frac12k_BT$）。由式 \eqref{eq:1d-E-avg}：
\begin{equation}
\frac12 mA^2\omega^2 = k_BT_m.
\end{equation}

题设熔点时振动振幅为平衡间距的 $10\%$，即
\begin{equation}
A = 0.1\,a = \frac{a}{10}.
\end{equation}
代入解出 $\omega$：
\begin{equation}
\omega = \frac{1}{A}\sqrt{\frac{2k_BT_m}{m}} = \frac{10}{a}\sqrt{\frac{2k_BT_m}{m}}.
\end{equation}
证毕。
\end{derivationbox}

\begin{formulabox}[答案汇总]
\begin{center}
\begin{tabular}{cl}
\toprule
问 & 结果 \\
\midrule
(1) 时间平均总能量 & $\displaystyle \langle E\rangle = \frac12 m\omega^2 A^2$ \\[8pt]
(2) 熔点附近振动频率 & $\displaystyle \omega = \frac{10}{a}\sqrt{\frac{2k_BT_m}{m}}$ \\
\bottomrule
\end{tabular}
\end{center}
\end{formulabox}

\begin{insightbox}[物理图像与深层联系]
\textbf{1. 为什么 $\langle K\rangle=\langle U\rangle$？——能量均分定理与``自由度''的准确含义}

对任何简谐振子，动能正比于 $\sin^2(\omega t+\phi)$，势能正比于 $\cos^2(\omega t+\phi)$，时间平均必然相等。这背后更根本的原理是\textbf{经典能量均分定理}：

> 哈密顿量中每一个\textbf{独立平方项}的平均能量为 $\frac12 k_BT$。

一维谐振子的哈密顿量 $H=\frac{p^2}{2m}+\frac12 m\omega^2 x^2$ 含有\textbf{两个}独立平方项（动量 $p$ 和坐标 $x$ 各一个），因此：
\begin{itemize}
    \item 动能项 $\frac{p^2}{2m}$ → $\langle K\rangle = \frac12 k_BT$
    \item 势能项 $\frac12 m\omega^2 x^2$ → $\langle U\rangle = \frac12 k_BT$
    \item 总能量 $\langle E\rangle = k_BT$
\end{itemize}

\textit{为什么固体物理常说``3个自由度''却给出 $3k_BT$？} 因为这里``自由度''指空间维度（$x,y,z$），而\textbf{每个空间维度}都包含动能+势能两个平方项，共 6 项，每项 $\frac12 k_BT$，所以总能量 $3k_BT$。这就是杜隆-珀蒂定律的来源。

\textbf{2. 势能计算中的 $1/2$ 因子陷阱}

初学者常犯的错误：
\begin{itemize}
    \item 错误 1：直接用 $U_n=\frac12\beta(x_n-x_{n+1})^2$，忘了每根键由两个原子共享，多算了一倍。
    \item 错误 2：把所有与原子 $n$ 相连的键（左右两根）都算作 $U_n$ 的独享，得到 $U_n=\frac14\beta[(x_n-x_{n-1})^2+(x_n-x_{n+1})^2]$。这在数学上其实等价于正确结果（因为每根键只出现一次），但概念上容易混淆。
\end{itemize}
最清晰的归属方式是：\textbf{先算单键能量，再乘以 $1/2$ 归给一个原子}。

\textbf{3. 为什么熔点附近满足经典极限？}

经典极限的严格条件是 $k_BT \gg \hbar\omega$。以典型金属 Cu 为例：
\begin{itemize}
    \item 德拜温度 $\Theta_D \approx 343$ K，爱因斯坦温度 $\Theta_E\sim\Theta_D$
    \item 熔点 $T_m \approx 1356$ K
    \item $\displaystyle\frac{\hbar\omega_E}{k_BT_m} = \frac{\Theta_E}{T_m} \sim \frac{300}{1356} \approx 0.22 \ll 1$
\end{itemize}

\textit{更直观的判据}：平均声子数 $\langle n\rangle \approx k_BT/\hbar\omega \gg 1$。熔点附近原子振幅达到晶格常数的 $\sim10\%$，振动能量 $k_BT_m\sim 4\hbar\omega_E$，对应的量子数很大，能级分立性被``抹平''，量子效应不可分辨——这正是经典行为。

\textbf{4. Lindemann 判据的简化版}

本题是\textbf{Lindemann 熔化判据}的一维简化形式。该判据认为：当原子热振动的均方根振幅达到晶格常数的某一临界比例（三维通常取 $10\%\text{--}15\%$）时，晶格失稳熔化。本题直接给定临界比例 $\eta=A/a=0.1$，结合能量均分定理即可导出熔点频率。

从式 \eqref{eq:1d-E-avg} 出发，可反解熔点温度：
\[
T_m = \frac{m\omega^2 A^2}{2k_B} = \frac{m\omega^2 a^2\eta^2}{2k_B}.
\]
若已知晶格参数 $a$ 和力常数 $\beta$（从而确定 $\omega$），即可估算熔点。这是固体物理中从微观参数预测宏观相变温度的经典范例。

\textbf{4. 与德拜模型的联系}

在德拜模型中，晶格的总零点能和热容均由声子模式的能量叠加得到。单个模式的平均能量正是 $\langle E\rangle=\hbar\omega(n_B+1/2)$。本题的经典极限对应高温 $k_BT\gg\hbar\omega$，此时 $\langle E\rangle\approx k_BT$，与量子结果自然衔接。
\end{insightbox}

\begin{thinkbox}[考场速记：简谐链能量问题的三步走]
\begin{enumerate}
    \item \textbf{写动能}：$K=\frac12 m\dot x^2$，时间平均得 $\langle K\rangle=\frac14 m\omega^2 A^2$
    \item \textbf{写势能}：先算键能 $U_{\text{bond}}=\frac12\beta(\Delta x)^2$，单原子分摊乘 $1/2$
    \item \textbf{用色散化简}：把 $\beta\sin^2(qa/2)$ 换成 $\frac14 m\omega^2$，得到 $\langle U\rangle=\langle K\rangle$
\end{enumerate}
\textit{检验}：若结果中 $\langle K\rangle\neq\langle U\rangle$，必定在某一步漏了因子 $2$ 或 $1/2$。
\end{thinkbox}
